{{دیوان
| عنوان = ہے یاد مجھے نکتۂ سلمانِ خوش آہنگ
| ماخذ = ویکی ماخذ
| ماخذ_لنک = https://ur.wikisource.org/wiki/%DB%81%DB%92_%DB%8C%D8%A7%D8%AF_%D9%85%D8%AC%DA%BE%DB%92_%D9%86%DA%A9%D8%AA%DB%82_%D8%B3%D9%84%D9%85%D8%A7%D9%86%D9%90_%D8%AE%D9%88%D8%B4_%D8%A2%DB%81%D9%86%DA%AF
| مقابلہ = اقبال سائبر لائبریری (اقبال اکادمی پاکستان)
| مقابلہ_لنک = https://www.iqbalcyberlibrary.net
}}
<poem>
ہے یاد مجھے نکتۂ سلمانِ خوش آہنگ
دنیا نہیں مردانِ جفاکش کے لیے تنگ

چیتے کا جگر چاہیے، شاہیں کا تجسّس
جی سکتے ہیں بے روشنیِ دانش و فرہنگ

کر بُلبل و طاؤس کی تقلید سے توبہ
بُلبل فقط آواز ہے، طاؤس فقط رنگ!
</poem>

٭ سلمان: مسعود سور سلیمانؔ۔ غزنوی دور کانامور ایرانی شاعر جو غالبا لاہور میں پیدا ہوا
